\section{Instruments 最有潜力候选总览}

\subsection{为什么单独抽这一组}

如果只看当前最值得继续讲、继续扩、继续做 follow-up 的 Instruments 候选，最有代表性的不是全部 family，而是下面这几条：

\begin{itemize}
\item \texttt{SFT495}：所有 RL 线的共同参考点。
\item \texttt{rule-only rerun}：当前最强的无 hint top-10 baseline。
\item \texttt{single-hint mixed}：当前最值得认真对待的新主候选。
\item old \texttt{fixed mixed-single}：带 bug caveat 的历史上界。
\item corrected \texttt{fixed taskfix sid-only}：当前最干净的 fixed reference。
\item corrected \texttt{fixed taskfix}：coverage 很强的 clean fixed 变体。
\item \texttt{hintce-3}：当前最值得继续追的 CE scaling 候选。
\end{itemize}

这一节不再把四个指标塞进一张 \texttt{2x2} 大图，而是每个指标单独出一张图。这样做的目的很直接：这几条线已经非常接近，拆成单指标之后更容易看清谁在 top-10 上领先，谁在 coverage 上更稳。

\subsection{最好点摘要}

\begin{table}[H]
\centering
\small
\setlength{\tabcolsep}{4pt}
\begin{tabular}{p{3.5cm} p{2.3cm} c c c c}
\toprule
Variant & Best ckpt & NDCG@10 & HR@10 & NDCG@50 & HR@50 \\
\midrule
SFT495 & \texttt{checkpoint-495} & 0.0823 & 0.1094 & 0.0985 & 0.1844 \\
rule-only rerun & \texttt{checkpoint-2997} & 0.0960 & 0.1179 & 0.1070 & 0.1681 \\
single-hint mixed & \texttt{checkpoint-2664} & 0.0948 & 0.1180 & 0.1116 & 0.1958 \\
old fixed mixed-single & \texttt{checkpoint-3326} & 0.0953 & 0.1193 & 0.1114 & 0.1938 \\
fixed taskfix sid-only & \texttt{checkpoint-2652} & 0.0945 & 0.1205 & 0.1103 & 0.1935 \\
fixed taskfix & \texttt{checkpoint-2997} & 0.0931 & 0.1189 & 0.1094 & 0.1941 \\
fixed CE coef=0.005 & \texttt{checkpoint-1665} & 0.0945 & 0.1180 & 0.1118 & 0.1985 \\
\bottomrule
\end{tabular}
\caption{Instruments 当前最有潜力候选的最好点。}
\end{table}

这张表最值得记住的不是单一冠军，而是当前候选的形状：

\begin{itemize}
\item \texttt{rule-only} 仍然是 top-10 最强的极值线，但 \texttt{HR@50} 明显最差。
\item \texttt{single-hint mixed}、old fixed、clean fixed family 和 \texttt{hintce-3} 已经整体挤进了更右上的一团，因此这几条线才是真正应该继续互相比的“前沿集合”。
\item \texttt{hintce-3} 当前 best 点的 \texttt{HR@50=0.1985} 最强，但它的 best 点更靠中段；\texttt{single-hint mixed} 和 old fixed 则更像完整长跑后的高位平台。
\end{itemize}

\subsection{NDCG@10}

\begin{figure}[H]
\centering
\includegraphics[width=0.96\textwidth]{promising_candidates_ndcg10.png}
\caption{Instruments promising candidates: NDCG@10。}
\end{figure}

\subsection{HR@10}

\begin{figure}[H]
\centering
\includegraphics[width=0.96\textwidth]{promising_candidates_hr10.png}
\caption{Instruments promising candidates: HR@10。}
\end{figure}

\subsection{NDCG@50}

\begin{figure}[H]
\centering
\includegraphics[width=0.96\textwidth]{promising_candidates_ndcg50.png}
\caption{Instruments promising candidates: NDCG@50。}
\end{figure}

\subsection{HR@50}

\begin{figure}[H]
\centering
\includegraphics[width=0.96\textwidth]{promising_candidates_hr50.png}
\caption{Instruments promising candidates: HR@50。}
\end{figure}

\subsection{怎么读这一组图}

\begin{itemize}
\item 如果你只关心 headline top-10，\texttt{rule-only} 仍然是最靠上的线，但它在 \texttt{HR@50} 图里会立刻暴露 coverage 损失。
\item 如果你关心 clean fixed 主线，\texttt{fixed taskfix sid-only} 和 \texttt{fixed taskfix} 现在已经足够稳定，而且一个更偏 top-10，一个更偏 coverage。
\item 如果你关心新的最强候选，当前真正需要继续盯的不是单条 fixed，而是 \texttt{single-hint mixed} 和 \texttt{hintce-3}：
  前者更像完整长跑后的主候选，后者更像中段爆发出的高 coverage / 高 NDCG 新候选。
\item old \texttt{fixed mixed-single} 仍然有必要保留在图里，因为它提供了历史上界；但它只能当 historical reference，不能再当 clean 方法定义。
\end{itemize}
